\chapter{Protocolo SDD/TDD e Pipeline de Qualidade para Agentes Cognitivos}

\begin{quote}
\textit{``Se você não pode medi-lo, não pode melhorá-lo.''} --- Lord Kelvin
\end{quote}

\section{Specification-Driven Development (SDD)}

\subsection{O Princípio: Especificar Antes de Implementar}

\textbf{SDD} inverte a prática comum de ``codar primeiro, documentar depois''. Toda tarefa começa com uma \textbf{especificação formal} (\texttt{specs/SPEC-*.md}) que define:

\begin{itemize}
    \item \textbf{Critérios de aceitação}: condições necessárias e suficientes para considerar a entrega completa;
    \item \textbf{Interface}: assinatura de funções, schemas de entrada/saída, contratos;
    \item \textbf{Restrições}: limites de performance, segurança, conformidade;
    \item \textbf{Testes}: cenários de teste associados (TDD).
\end{itemize}

\begin{lstlisting}[language=Python, caption={SpecRegistry — Gerenciador de Especificações Formais}]
class SpecRegistry:
    """Registro central de especificações SDD."""
    
    def __init__(self, specs_dir="specs/"):
        self.specs = {}
        self._load_all(specs_dir)
    
    def _load_all(self, specs_dir):
        import glob
        for spec_file in glob.glob(f"{specs_dir}/SPEC-*.md"):
            spec = Specification.from_markdown(spec_file)
            self.specs[spec.spec_id] = spec
    
    def get(self, spec_id):
        if spec_id not in self.specs:
            raise SpecNotFoundError(f"SPEC {spec_id} não encontrada")
        return self.specs[spec_id]
    
    def verify(self, spec_id, deliverable):
        """Verifica se a entrega satisfaz a especificação."""
        spec = self.get(spec_id)
        results = []
        for criterion in spec.acceptance_criteria:
            passed = criterion.evaluate(deliverable)
            results.append({'criterion': criterion.description, 'passed': passed})
        return all(r['passed'] for r in results), results
\end{lstlisting}

\section{Test-Driven Development (TDD) para Agentes}

\subsection{O Ciclo RED → GREEN → REFACTOR}

Beck (2003) \cite{beck2003tdd} estabeleceu o ciclo TDD que o OpenCode aplica rigorosamente:

\begin{enumerate}
    \item \textbf{RED}: escrever um teste que falha (porque a funcionalidade ainda não existe);
    \item \textbf{GREEN}: implementar o código mínimo para o teste passar;
    \item \textbf{REFACTOR}: melhorar o código mantendo os testes verdes.
\end{enumerate}

Para agentes, isto se traduz em:
\begin{itemize}
    \item \textbf{Testes de comportamento}: verificar outputs esperados para inputs conhecidos;
    \item \textbf{Testes de robustez}: inputs adversariais, edge cases, condições de erro;
    \item \textbf{Testes de metacognição}: verificar se o agente reflete sobre falhas e melhora.
\end{itemize}

\subsection{PRÁXIS 8.1 — TDD para um Agente com pytest}

\begin{lstlisting}[language=Python, caption={TDD: teste primeiro, implemente depois}]
# tests/test_agent_coder.py
import pytest
from agents.coder import CoderAgent

class TestCoderAgent:
    """Bateria TDD para o agente Coder."""
    
    def setup_method(self):
        self.agent = CoderAgent(spec_id="SPEC-006")
    
    def test_red_generate_function_stub(self):
        """RED: função ainda não existe."""
        with pytest.raises(NotImplementedError):
            self.agent.generate("sort array", language="python")
    
    def test_green_implement_sort(self):
        """GREEN: implementação mínima."""
        self.agent.implement_sort()
        code = self.agent.generate("sort array", language="python")
        assert "def sort" in code
        assert "return" in code
    
    def test_refactor_optimize(self):
        """REFACTOR: otimização sem quebrar testes."""
        code_v1 = self.agent.generate("sort array", language="python")
        self.agent.optimize()
        code_v2 = self.agent.generate("sort array", language="python")
        # Mesma funcionalidade, mas mais eficiente
        assert len(code_v2) <= len(code_v1)
        assert "sorted(" in code_v2 or ".sort()" in code_v2
\end{lstlisting}

\section{Gate SDD + BehavioralGate: Dupla Verificação}

\subsection{Arquitetura de Duplo Gate}

O pipeline de qualidade do OpenCode implementa \textbf{dois gates sequenciais}:

\begin{enumerate}
    \item \textbf{Gate SDD (SpecVerifier)}: verifica se a entrega satisfaz a especificação formal (critérios de aceitação). É um gate \textbf{determinístico}: passa/falha binário.
    
    \item \textbf{BehavioralGate (Trust Engine)}: avalia a qualidade comportamental do agente ao longo do tempo. É um gate \textbf{probabilístico}: calcula Trust Score via EMA e decide com base em threshold.
\end{enumerate}

Uma entrega só é aceita se passar em \textbf{ambos} os gates:

\begin{equation}
\text{Accept}(\text{deliverable}) = \text{SpecVerifier}(d) \land (\text{TrustScore}(a) \geq \theta_{\text{trust}})
\label{eq:double_gate}
\end{equation}

Onde $\theta_{\text{trust}} = 0.7$ é o threshold padrão (agentes com Trust Score $< 0.3$ exigem supervisão obrigatória).

\section{Trust Engine e Confidence Ledger}

\subsection{Mecanismo de Confiança Adaptativa}

O \textbf{Trust Engine} (\texttt{trust/trust\_engine.py}) calcula e mantém o Trust Score de cada agente:

\begin{equation}
S_{t+1}(a) = \alpha \cdot \text{outcome}_t + (1-\alpha) \cdot S_t(a)
\label{eq:trust_ema}
\end{equation}

Onde $\alpha = 0.3$ (30\% de peso para o resultado mais recente, 70\% para o histórico). Este é essencialmente um \textbf{filtro de média móvel exponencial} (EMA), comum em finanças e processamento de sinais.

\subsection{Natural Forgetting e Decaimento}

Agentes inativos sofrem \textbf{decaimento natural de confiança}:

\begin{equation}
S_{t+\Delta}(a) = S_t(a) \cdot \gamma^{\Delta}
\label{eq:trust_decay}
\end{equation}

Onde $\gamma \in (0, 1)$ é o fator de decaimento por unidade de tempo $\Delta$. Isso evita que agentes mantenham scores altos após longos períodos sem atividade.

\section{Slashing e Incentivos na Token Economy}

\subsection{Curva de Penalidade Proporcional}

O \textbf{slashing} é a penalidade econômica por falha, proporcional ao dano causado:

\begin{equation}
\text{slash}(a, \text{failure}) = \text{stake}(a) \cdot \min\left(\text{severity}(\text{failure}), \text{max\_slash\_ratio}\right)
\label{eq:slashing}
\end{equation}

A severidade é calculada pelo BehavioralGate:
\begin{itemize}
    \item \textbf{Nível 1} (leve): output abaixo do esperado mas funcional — slash 5\%;
    \item \textbf{Nível 2} (moderado): output incorreto ou incompleto — slash 25\%;
    \item \textbf{Nível 3} (grave): output danoso, violação ética ou falha de segurança — slash 100\%.
\end{itemize}

\begin{thebibliography}{99}
\bibitem{beck2003tdd} Beck, K. (2003). \textit{Test-Driven Development: By Example}. Addison-Wesley.
\bibitem{meyer1997object} Meyer, B. (1997). \textit{Object-Oriented Software Construction} (2nd ed.). Prentice Hall.
\bibitem{fowler2018refactoring} Fowler, M. (2018). \textit{Refactoring: Improving the Design of Existing Code} (2nd ed.). Addison-Wesley.
\bibitem{shinn2023reflexion} Shinn, N., et al. (2023). Reflexion. \textit{NeurIPS 2023}. \doiurl{10.48550/arXiv.2303.11366}
\bibitem{nash1950equilibrium} Nash, J. F. (1950). Equilibrium Points in N-Person Games. \textit{PNAS}, 36(1), 48--49. \doiurl{10.1073/pnas.36.1.48}
\bibitem{buterin2014ethereum} Buterin, V. (2014). Ethereum: A Next-Generation Smart Contract Platform.
\bibitem{voshmgir2020token} Voshmgir, S. (2020). \textit{Token Economy}. Token Kitchen.
\bibitem{myerson1981optimal} Myerson, R. B. (1981). Optimal Auction Design. \textit{Mathematics of Operations Research}, 6(1), 58--73. \doiurl{10.1287/moor.6.1.58}
\end{thebibliography}
